\section{RQ4: Trends}
\label{sec:trends}

% Included studies per 1,000 accepted papers at venue-years fully covered by Source B.
\begin{figure}[t]
\centering
\begin{tikzpicture}
\begin{axis}[
  width=\columnwidth, height=50mm,
  xmin=2019.6, xmax=2026.4, ymin=0,
  xtick={2020,2021,2022,2023,2024,2025,2026},
  xticklabel style={font=\scriptsize, /pgf/number format/1000 sep={}},
  yticklabel style={font=\scriptsize},
  ylabel={Included per 1{,}000 accepted}, ylabel style={font=\scriptsize},
  axis lines*=left, ymajorgrids, grid style={black!10},
  legend style={font=\tiny, at={(0.02,0.98)}, anchor=north west, draw=none, fill=white, fill opacity=0.85, text opacity=1},
  legend cell align=left,
]
\addplot[nsblue, line width=1pt, mark=*, mark size=1.6pt]
  table[col sep=comma, x=year, y=panel5_per1000, restrict x to domain=2020:2025] {data/derived/plot_growth.csv};
\addplot[nsgrey, line width=0.8pt, dashed, mark=square*, mark size=1.4pt]
  table[col sep=comma, x=year, y=panel2_per1000] {data/derived/plot_growth.csv};
\legend{ICLR, ICML, NeurIPS, ACL, EMNLP (2020--25), ICLR and ICML (2020--26)}
\end{axis}
\end{tikzpicture}
\caption{Growth normalised by the number of accepted papers at the same venue-years (Source B). The five-venue panel rises from \cnt{pfive2020} per 1{,}000 accepted papers in 2020 to \cnt{pfive2025} in 2025; the ICLR+ICML panel reaches \cnt{ptwo2026} in 2026. Only \cnt{nesytitle} of the \cnt{included} included studies (\cnt{nesytitlepct}\%) use ``neurosymbolic'' or ``neural-symbolic'' in their title, so the trend is not driven by adoption of the term.}
\label{fig:trend}
\end{figure}


\textbf{Growth that is not a terminology artefact.} The pilot reported growth as the share of a title-keyword corpus, which would also rise if recent authors were simply more likely to call their work ``neurosymbolic''. Two measures address this (Fig.~\ref{fig:trend}).

First, growth is normalised by all accepted papers at venue-years with constant coverage. In the five-venue panel (ICLR, ICML, NeurIPS, ACL, EMNLP) included studies rise from \cnt{pfive2020} per 1{,}000 accepted papers in 2020 to \cnt{pfive2023} in 2023 and \cnt{pfive2025} in 2025. The denominators are \cnt{accfive2020} and \cnt{accfive2025} accepted papers. In the ICLR+ICML panel the rate reaches \cnt{ptwo2026} in 2026.

Second, inclusion no longer depends on the term: only \cnt{nesytitle} included studies (\cnt{nesytitlepct}\%) use ``neurosymbolic'' or ``neural-symbolic'' in the title, and \cnt{nesyany} (\cnt{nesyanypct}\%) anywhere in title or abstract. The rise is therefore real growth of the area, a roughly \cnt{growthratio}-fold increase in relative share over five years. It is not adoption of a word.

Two caveats apply. Search recall may differ by year if vocabulary changed; the audit is too small to estimate recall per year. The 2026 figures cover only the venues held by October 2026.

% Integration architectures per period (data/derived/arch_by_period.csv), stacked counts.
\begin{figure}[t]
\centering
\pgfplotstableread[col sep=comma]{data/derived/arch_by_period.csv}\archperiodtable
\begin{tikzpicture}
\begin{axis}[
  width=\columnwidth, height=50mm,
  xbar stacked, bar width=9pt,
  xmin=0, xlabel={Included studies}, xlabel style={font=\scriptsize},
  ytick=data, yticklabels from table={\archperiodtable}{period},
  y dir=reverse,
  yticklabel style={font=\scriptsize}, xticklabel style={font=\scriptsize},
  legend style={font=\tiny, at={(0.5,1.03)}, anchor=south, legend columns=7, draw=none, /tikz/every even column/.append style={column sep=2pt}},
  legend image code/.code={\draw[#1, draw=none] (0cm,-0.08cm) rectangle (0.22cm,0.12cm);},
  axis lines*=left, xmajorgrids, grid style={black!10},
  enlarge y limits=0.2,
]
\addplot[fill=A1c] table[col sep=comma, x=A1, y expr=\coordindex] {data/derived/arch_by_period.csv};
\addplot[fill=A2c] table[col sep=comma, x=A2, y expr=\coordindex] {data/derived/arch_by_period.csv};
\addplot[fill=A3c] table[col sep=comma, x=A3, y expr=\coordindex] {data/derived/arch_by_period.csv};
\addplot[fill=A4c] table[col sep=comma, x=A4, y expr=\coordindex] {data/derived/arch_by_period.csv};
\addplot[fill=A5c] table[col sep=comma, x=A5, y expr=\coordindex] {data/derived/arch_by_period.csv};
\addplot[fill=A6c] table[col sep=comma, x=A6, y expr=\coordindex] {data/derived/arch_by_period.csv};
\addplot[fill=A7c] table[col sep=comma, x=A7, y expr=\coordindex] {data/derived/arch_by_period.csv};
\legend{A1,A2,A3,A4,A5,A6,A7}
\end{axis}
\end{tikzpicture}
\caption{Integration architectures by period. Before 2024, symbolic guidance (A3) dominates; from 2025, translate-and-solve (A1) and verifier-in-the-loop (A2) designs together account for \cnt{y2025AoneAtwo} of \cnt{y2025total} studies in 2025 and \cnt{y2026AoneAtwo} of \cnt{y2026total} in 2026, against \cnt{earlyAoneAtwo} of \cnt{earlytotal} before 2024. $^\ast$Partial year.}
\label{fig:archperiod}
\end{figure}


\textbf{From guidance to checking.} The architecture mix shifted (Fig.~\ref{fig:archperiod}).
\begin{itemize}
  \item \emph{2020--2022} ($n=\cnt{earlyn}$; with 2023 the pre-2024 sample has \cnt{pretwentyfour} studies, against 8 in the pilot). Symbolic guidance (A3) accounts for \cnt{earlyAthree}\% of studies: constrained decoding, semantic parsing and logical decoding constraints~\cite{scholak2021picard,lu2021neurologic,shin2021constrained}.
  \item \emph{2026.} Checker-in-the-loop designs (A2) account for \cnt{ysixAtwo}\% [95\% CI \cnt{ysixAtwolo}--\cnt{ysixAtwohi}], up from \cnt{earlyAtwo}\% [\cnt{earlyAtwolo}--\cnt{earlyAtwohi}]; A3 falls to \cnt{ysixAthree}\% [\cnt{ysixAthreelo}--\cnt{ysixAthreehi}]. The intervals for the first and last periods do not overlap.
  \item \emph{Exact checks.} The share of studies whose symbolic side gives an exact verdict rises from \cnt{earlyexact}\% to \cnt{ysixexact}\% [\cnt{ysixexactlo}--\cnt{ysixexacthi}], driven by formal theorem proving and verified code generation.
\end{itemize}
Checking designs need a model capable of emitting well-formed formal artifacts, which is consistent with their timing. Translate-and-solve designs (A1) hold about a fifth of studies until 2025 (\cnt{yfiveAone}\%) and fall to \cnt{ysixAone}\% in 2026, consistent with single-pass formalisation giving way to checked loops.

\textbf{Two communities that rarely meet.} Classical neurosymbolic themes rarely appear in the included set. Probabilistic logic programming, abductive learning, inductive logic programming and reasoning-shortcut analysis are well represented in the curated corpus~\cite{huang2021scallop,maene2024hardness,marconato2023neuro,yang2024analysis}. Yet few included studies use these semantics or diagnostics with language models; examples are neurosymbolic losses and tractable probabilistic models for constrained generation~\cite{ahmed2023pseudo,ahmed2025controllable,zhang2023tractable}. The language-model line has not yet adopted the classical line's tools for calibrated inference and for detecting satisfied-for-the-wrong-reason symbols.
